\documentclass[a4paper]{article}

\usepackage{graphicx}
\usepackage{amsmath,amssymb}
\usepackage{minted}
\usepackage{multirow}
\usepackage{float}
\usepackage{todonotes}
\usepackage{forest}
\usepackage{xstring}
\usepackage{xcolor}
\usepackage[most]{tcolorbox}
\usepackage{xparse}
\usepackage{natbib}

\usetikzlibrary{graphs,graphdrawing}
\usegdlibrary{layered}

% for external graphviz
\input{/home/donkarlo/Dropbox/repo/nd_ai_project/src/nd_ai/knoweldge_representation/language/formal/computer/programming/tex/latex/code_clips/external_graphviz.tex}

% for tags
\input{/home/donkarlo/Dropbox/repo/nd_ai_project/src/nd_ai/knoweldge_representation/language/formal/computer/programming/tex/latex/parts/control_sequence/kind/semtag/semtag.tex}

% for links
\input{/home/donkarlo/Dropbox/repo/nd_ai_project/src/nd_ai/knoweldge_representation/language/formal/computer/programming/tex/latex/code_clips/seehere.tex}

\title{Bei}
\date{}
\author{Mohammad Rahmani}

\begin{document}
    \maketitle
    
    
    \section{Kasus}
        \begin{itemize}
            \item Bei ist immer mit dativ
        \end{itemize}
    
    
    \section{Die Anwendungsfälle von bedeutung Perspektiv}
        
        \subsection{beim letzten}
            beim letzten: at the last one; during the last; on the last occasion
            
            Bedeutung: Kurzform von „bei dem letzten“. Es bedeutet je nach Kontext: bei der letzten Person/Sache, beim letzten Mal oder während des letzten Ereignisses.
            
            ex: Beim letzten Mal war ich krank: Last time I was sick.
            
            ex: Beim letzten Versuch hat es funktioniert: It worked on the last attempt.
            
            ex: Beim letzten Termin war er nicht da: He was not there at the last appointment.
            
            
            \bibliographystyle{plainnat}
            \bibliography{/home/donkarlo/Dropbox/repo/nd_shared_research/bibliography/bibliography.bib}
\end{document}